% Physical unit 027; accepted Session 03 after fresh 400-dpi authority replay.
\GordanUnitStart{027}{printed page 22}
\GordanTopFolio{22}
\begingroup\small
the point $C$. Since the length of each of these two lines between $A$ and $C$ is the same, there exist from $A$ to $B$ two equal paths, one of which is geodetic line I itself, the other of which is made up of two parts, namely geodetic line II between $A$ and $C$ and geodetic line I between $C$ and $B$.

But if the former path between $A$ and $B$ were the shortest of all lines, then certainly the latter also would be the shortest of all lines, which cannot happen, because the latter path is not a geodetic line at the point $C$.

II. ``That geodetic line $l_0$ of the first or third genus which, drawn from $A$ to $B$, cuts the circle $-l_0$ twice (one of the places indeed being in $C$), is not the shortest of all lines between $A$ and $B$.''

The reason for demonstrating this is similar to the preceding one.
For with the point $\Pi$ carried from $A$ on the other geodetic line, which touches the circle $l_0$, to $C$, and then on the former from $C$ to $B$, this whole path is equal to the length $AB$ of geodetic line I.

Thus what we asserted has been demonstrated by the preceding mode of argument.

\begin{center}\S\ 6.\end{center}

The line touching the circle $l_1$ at the point $A$ divides the zone enclosed by the circles $l_1$ and $-l_1$ into two parts, one of which we choose to name ``$J$'', lying inside the geodetic line, the other ``$E$'', lying outside the geodetic line.

Therefore, once the force contained in the theorems just proposed has been recognized, it is clear that the shortest of all lines between $A$ and $B(l\psi)$, when $B$ lies in the part ``$J$'', is of the first or third genus and has no point of conversion between $A$ and $B$; when $B$ lies in the part ``$E$'', it is of the second or fourth genus and has a point of conversion between $A$ and $B$; if
\[
 \psi>\omega,
\]
it is of the first or second genus, while if
\[
 \psi<\omega,
\]
it is of the third or fourth genus.

Thus, whatever part of the zone contains the point $B$, it can be determined in which genus of geodetic lines the shortest of all lines between $A$ and $B$ is to be placed; and indeed the geodetic line which is shortest of all does not cut the circle $-l_0$.

\begin{center}\S\ 7.\end{center}

These facts tend to this: if a geodetic line is part of the length of a period, it is shorter than the other geodetic lines between $A$ and $B$; or,
``The geodetic line $l_0$ is, under this condition, the shortest of all lines between $A$ and $B$, if the difference of the angles $\varphi$ belonging to these points does not exceed $180^\circ$.''
\endgroup
